% !PRINTED 127
\weberpage{127}
%<<CONT
is a power of $p$. $R'$ is the totality of all elements of $P$ permutable with $Q'$.

We now further apply the theorem 3., \S.\,28, understanding by $P$, $R$, $Q$ the groups that are here denoted by the same letters. Since here $Q$ is a divisor of $R$, we have $h_1 = 1$. But of the conjugate groups $b^{-1}Rb$ none except $R$ is divisible by $Q$, and so the remaining $h_2, h_3, \ldots$ are all greater than $1$. But since they must be divisors of the degree of $Q$, they are powers of $p$, and from the equation

\[
j = h_1 + h_2 + h_3 + \cdots
\]

it therefore follows that $j \equiv 1\ (\mathrm{mod}\ p)$.

But if $Q_1$ is any divisor of $P$ whose degree is a power of $p$, we again apply the decomposition of $j$ according to the Theorem 3., \S.\,28, putting for the groups $P$, $R$, $Q$ occurring there the groups $P$, $R$, $Q_1$. Then again

\[
j = h_1 + h_2 + h_3 + \cdots,
\]

in which now $h_1, h_2, h_3, \ldots$ denote the indices (with respect to $Q$) of the intersections of $Q_1$ with the group $R$ and its conjugates, which are therefore likewise powers of $p$. % [sic: Weber prints "in Bezug auf Q"; after substituting Q_1 for theorem 3's Q, these indices are with respect to Q_1.]
But since $j \equiv 1\ (\mathrm{mod}\ p)$, for at least one of the numbers we must have $h = 1$, i.e. one of the groups $b_\lambda^{-1}Rb_\lambda$, and consequently also one of the groups (11), is divisible by $Q_1$. If $Q_1$ is of degree $p^\varkappa$, it is identical with one of the groups (11).

With this, then, the Theorem II is completely proved.

\begin{center}
\S.\,31.

\gesperrt{Groups of degree $p^\alpha$.}
\end{center}

\gesperrt{Sylow} has derived from his theorems a remarkable property of the groups whose degree is a power of a prime $p$, namely $p^\alpha$. Let $P$, then, be such a group and $\beta$ a positive number smaller than $\alpha$. Then by the theorem I, \S.\,29 the group $P$ contains a divisor $Q$ of degree $p^\beta$, whose index is $p^{\alpha-\beta}$. We apply the theorem 3., \S.\,28, putting for $R$ and $Q$ this group $Q$ of degree $p^\beta$; then $h_1, h_2, h_3, \ldots$ are
%CONT>>
